% SEGMENT OLP-0597-S01
\documentclass[../../../include/open-logic-section]{subfiles}

\begin{document}

\olfileid{sth}{choice}{countablechoice}
\olsection{எண்ணத்தக்க தேர்வு}

தேர்வையோ நன்கு வரிசைப்படுத்தலையோ பயன்படுத்தாமல்
பின்வருவதை எளிதில் நிறுவலாம்:

\begin{lem}[$\Zminus$ இல்]
ஒவ்வொரு முடிவுறு கணத்துக்கும் ஒரு தேர்வுச் சார்பு உள்ளது.
\end{lem}

\begin{proof}
$a = \{b_1, \ldots, b_n\}$ என்க. எளிமைக்காக ஒவ்வொரு
$b_i \neq \emptyset$ என்றும் கொள்வோம். அப்போது
$c_1 \in b_1, \ldots, c_n \in b_n$ ஆகும்
$c_1, \ldots, c_n$ என்ற பொருள்கள் உள்ளன.
\olref[z][pairs]{prop:pairsconsequences} இன்படி
$\{\langle b_1, c_1\rangle , \ldots, \langle b_n, c_n\rangle\}$
என்ற கணம் உள்ளது; இதுவே $a$ க்கான தேர்வுச் சார்பு.
\end{proof}

ஆனால் முடிவுறாத கணங்களைக் கருதியவுடன் நிலைமை
தெளிவற்றதாகிறது. மேலுள்ள முடிவின் மிகச் சிறிய நீட்டிப்பைப்
பார்க்கலாம்:

\begin{defish}
\emph{எண்ணத்தக்க தேர்வு.} ஒவ்வொரு \emph{எண்ணத்தக்க}
கணத்துக்கும் ஒரு தேர்வுச் சார்பு உள்ளது.
\end{defish}

இது தேர்வின் அடிகோளத்தின் ஒரு தனிநிகழ்வு. இந்தக் கொள்கை
பயன்படுத்தப்படுகிறது என்பதே வெளிப்படையாக உணரப்படாமல்,
பலமுறை பயன்படுத்தப்பட்டதும் தெரியவருகிறது. இரண்டு
எடுத்துக்காட்டுகளைப் பார்ப்போம்.\footnote{இவை
\citet[\S9.4]{Potter2004} மற்றும் லூகா இன்குர்வாட்டி
ஆகியோரிடமிருந்து பெறப்பட்டவை.}

\begin{ex}
இயல்பாகத் தோன்றும் ஓர் எண்ணம்: எந்தக் கணம் $A$ க்கும்,
$\cardle{\omega}{A}$ அல்லது ஏதேனும் $n \in \omega$
க்கு $\cardeq{A}{n}$. ஒவ்வொரு கணமும் முடிவுறு அல்லது
முடிவுறாதது என்ற உள்ளுணர்வுக் கருத்தைச் சொல்லும் ஒரு வழி
இது. காண்டோரும் பல கணிதவியலாளர்களும் இதை நிறுவாமல்
கூறினர். எச்சரிக்கையுடன் செயல்பட்ட நாம்,
\olref[cardinals][classing]{generalinfinitycharacter} இல்
இதைக் நிறுவினோம். ஆனால் அந்த நிறுவலில் எந்தக் கணம்
$A$ ஐயும் நன்கு வரிசைப்படுத்தலாம் என்றும், அதனால்
$\card{A}$ இருப்பது உறுதி என்றும் எடுத்துக்கொண்டதால்,
$\ZFC$ இல் வேலை செய்தோம். அதாவது தேர்வை வெளிப்படையாக
எடுத்துக்கொண்டோம்.

% SEGMENT OLP-0597-S02
உண்மையில் \citet{Dedekind1888} இந்தக் கூற்றுக்குத்
தனது நிறுவலைப் பின்வருமாறு வழங்கினார்:

\begin{thm}[$\Zminus + \text{எண்ணத்தக்க தேர்வு}$ இல்]
எந்த $A$ க்கும், $\cardle{\omega}{A}$ அல்லது ஏதேனும்
$n \in \omega$ க்கு $\cardeq{A}{n}$.
\end{thm}

\begin{proof}
எல்லா $n \in \omega$ க்கும் $\cardneq{A}{n}$ என்க.
அப்போது குறிப்பாக ஒவ்வொரு $n < \omega$ க்கும்,
சரியாக $\cardexpo{2}{n}$ உறுப்புகளைக் கொண்ட
$A_n \subseteq A$ என்ற துணைக்கணம் உள்ளது.
இந்த $A_0, A_1, A_2, \ldots$ தொடரைப் பயன்படுத்தி,
ஒவ்வொரு $n$ க்கும் பின்வருமாறு வரையறுக்கிறோம்:
\[
	B_n = A_n \setminus \bigcup_{i < n} A_i.
\]
இப்போது பின்வருவதை கவனிக்கவும்:
% SOURCE CORRECTION TA-STH-045: ஒன்றிப்பின் குறியீட்டை A_i எனச் சரிசெய்க.
\begin{align*}
	\card{\bigcup_{i < n}A_i}
	&\leq \card{A_0} + \card{A_1} + \ldots + \card{A_{n-1}}\\
	&=1 + 2 + \ldots + 2^{n-1}\\
	& = 2^n - 1\\
	& < 2^n = \card{A_n}
\end{align*}
ஆகவே ஒவ்வொரு $B_n$ க்கும் குறைந்தது ஒரு உறுப்பு $c_n$
உள்ளது. மேலும் $B_n$ கள் ஒன்றையொன்று வெட்டாதவை.
எனவே $c_n = c_m$ எனில் $n = m$. ஆனால் ஒவ்வொரு
$c_n \in A$; ஆகவே $f(n) = c_n$ என்ற சார்பு
$\omega \to A$ என்ற உள்செலுத்தல்.
\end{proof}
\noindent

% SEGMENT OLP-0597-S03
டெடெகிண்ட் இங்கு எண்ணத்தக்க தேர்வைப் பயன்படுத்தியதைக்
குறிப்பிடவில்லை. ஆனால் \emph{நீங்கள்} அதன் பயன்பாட்டைக்
கண்டீர்களா? மீண்டும் பாருங்கள். (உண்மையில்,
\emph{மீண்டும் பாருங்கள்}.)

இந்த நிறுவல் எண்ணத்தக்க தேர்வை இருமுறை பயன்படுத்தியது.
முதலில் $A_0$, $A_1$, $A_2$, \dots\@ என்ற கணங்களின்
தொடரைப் பெறப் பயன்படுத்தினோம். பின்னர் ஒவ்வொரு~$B_n$
இலிருந்தும் $c_n$ என்ற உறுப்பைத் தேர்ந்தெடுக்கப்
பயன்படுத்தினோம். மேலும், இந்தத் தேர்வைப் பயன்படுத்துவதை
நீக்க முடியாது. தேர்வின் எந்த வடிவமும் இல்லாவிட்டால்
முடிவு தவறக்கூடும் என்று \citet[p.~138]{Cohen1966}
நிறுவினார். அதாவது, $\omega$ வுடன் \emph{ஒப்பிட முடியாத}
கணங்கள் இருப்பது $\ZF$ உடன் முரண்படாது.
\end{ex}

% SEGMENT OLP-0597-S04
\begin{ex}
\citeyear{Cantor1878} இல், எண்ணத்தக்க கணங்களின்
எண்ணத்தக்க ஒன்றிப்பு எண்ணத்தக்கது என்று காண்டோர் கூறினார்.
நிறுவலை அவர் அளிக்கவில்லை; அது வெளிப்படையானது என்று
கருதியிருக்கலாம். நாம் \olref[card-arithmetic][simp]{kappaunionkappasize}
இல் இதைவிடப் பொதுவான முடிவை நிறுவினோம். ஆனால் நமது
நிறுவல் தேர்வை வெளிப்படையாக எடுத்துக்கொண்டது. குறைவான
பொதுத்தன்மை உடைய இந்த முடிவிற்குக்கூட எண்ணத்தக்க தேர்வு
தேவை.

\begin{thm}[$\Zminus + \text{எண்ணத்தக்க தேர்வு}$ இல்]
ஒவ்வொரு $n \in \omega$ க்கும் $A_n$ எண்ணத்தக்கது எனில்,
$\bigcup_{n < \omega} A_n$ எண்ணத்தக்கது.
\end{thm}

\begin{proof}
பொதுத்தன்மை இழப்பின்றி ஒவ்வொரு $A_n \neq \emptyset$
என்க. ஆகவே ஒவ்வொரு $n \in \omega$ க்கும்
$f_n \colon \omega \to A_n$ என்ற !!a{surjection}
உள்ளது. $f(m, n) = f_n(m)$ என வரையறுத்து,
$f \colon \omega \times \omega \to \bigcup_{n < \omega} A_n$
என்ற சார்பைப் பெறுகிறோம். $\omega \times \omega$
எண்ணத்தக்கது
(\olref[sfr][siz][zigzag]{natsquaredenumerable}) என்பதாலும்
$f$ ஒரு !!a{surjection} என்பதாலும் முடிவு கிடைக்கிறது.
\end{proof}
\noindent

% SEGMENT OLP-0597-S05
எண்ணத்தக்க தேர்வின் பயன்பாட்டை இங்கும் கண்டீர்களா?
$f_0$, $f_1$, $f_2$, \dots\footnote{இதைப் போன்ற தேர்வின்
பயன்பாடு \olref[card-arithmetic][simp]{kappaunionkappasize}
இலும் நடந்தது: ``ஒவ்வொரு $\beta \in \cardfont{a}$ க்கும்,
$f_\beta$ என்ற !!a{injection} ஐ நிலைப்படுத்துக'' என்று
அங்கே சொன்னோம்.} என்ற சார்புகளின் தொடரைத் தேர்ந்தெடுக்க
அது பயன்படுகிறது. தேர்வின் எந்தக் கொள்கையும் இல்லாதபோது
முடிவு மீண்டும் தவறலாம். குறிப்பாக, எண்ணத்தக்க கணங்களின்
எண்ணத்தக்க ஒன்றிப்பின் கண அளவு~$\beth_1$ ஆகும் நிலை
$\ZF$ உடன் முரண்படாது என்று
\citet{FefermanLevy1963} நிறுவினர். இதே கருத்தை
ரசல் சுவையாகச் சொல்கிறார்:
\begin{quote}
  ஒரு கோடீஸ்வரர் ஒவ்வொரு முறை ஒரு சோடி காலணிகளை
  வாங்கும்போதும் ஒரு சோடி காலுறைகளையும் வாங்கினார்;
  வேறு எந்த நேரத்திலும் வாங்கவில்லை என்று எண்ணுக.
  இரண்டையும் வாங்குவதில் அவருக்கு இருந்த ஆர்வத்தால்,
  இறுதியில் அவரிடம் $\aleph_0$ சோடி காலணிகளும்
  $\aleph_0$ சோடி காலுறைகளும் இருந்தன\dots\@ காலணிகளில்
  வலது, இடது என்பதை வேறுபடுத்தலாம்; ஆகவே ஒவ்வொரு
  சோடியிலிருந்தும் ஒன்றைத் தேர்ந்தெடுக்கலாம், எல்லா
  வலது காலணிகளையோ எல்லா இடது காலணிகளையோ எடுக்கலாம்.
  ஆனால் காலுறைகளுக்கு அத்தகைய தேர்வு விதி எதுவும்
  தோன்றவில்லை. பெருக்கல் அடிகோளத்தை [நடைமுறையில்,
  தேர்வை] ஏற்றுக்கொள்ளாவிட்டால், ஒவ்வொரு சோடியிலிருந்தும்
  ஒரு காலுறையைக் கொண்ட ஒரு கணம் இருப்பதை உறுதியாகக்
  கூற முடியாது.
  \citep[p.~126]{Russell1919}
\end{quote}
வேறுவிதமாகச் சொன்னால், எண்ணத்தக்க அளவில் காலுறைச்
சோடிகள் இருந்தால், மொத்தக் காலுறைகளும் எண்ணத்தக்க
அளவிலேயே இருக்கும் என்பதை நிறுவத் தேர்வின் ஏதேனும்
ஒரு வடிவம் தேவை. உண்மையில், எண்ணத்தக்க தேர்வு அல்லது
அதற்குச் சமமான வேறு கொள்கை இல்லாவிட்டால்,
எண்ணத்தக்க கணங்களின் எண்ணத்தக்க ஒன்றிப்பு எண்ணத்தக்கதாக
இல்லாமல் போகலாம்.
\end{ex}

% SEGMENT OLP-0597-S06
எண்ணத்தக்க தேர்வைப் பயன்படுத்தியவர்களே அதை அதிகம்
உணராமல் மீண்டும் மீண்டும் பயன்படுத்தினர் என்பதே கருத்து.
தத்துவக் கேள்வி: எண்ணத்தக்க தேர்வை நாம் எப்படி
\emph{நியாயப்படுத்தலாம்}?

உள்ளுணர்வால் நியாயப்படுத்த முயல்பவர் ஒரு அதிவேகப் பணியை
முன்வைக்கலாம். முதல் தேர்வுக்கு $\nicefrac{1}{2}$ நிமிடம்,
இரண்டாவதற்கு $\nicefrac{1}{4}$ நிமிடம், \dots,
$n$ ஆம் தேர்வுக்கு $\nicefrac{1}{2^n}$ நிமிடம்,
\dots\@ என எடுத்துக்கொள்வோம். அப்போது $1$~நிமிடத்திற்குள்
ஓர் $\omega$-நீளத் தேர்வுத் தொடரை முடித்து, ஒரு தேர்வுச்
சார்பை வரையறுத்திருப்போம்.

ஆனால் உண்மையில் இத்தகைய சிந்தனைப் பரிசோதனை என்ன
சொல்ல முடியும்? முதலாவதாக, ``தேர்ந்தெடுத்தல்'' என்ற கருத்தை
மிகவும் நேர்ப்பொருளில் எடுக்கிறது. அடுத்ததாக, கணிதத்தை
மெய்யியல் சாத்தியத்துடன் பிணைப்பதாகத் தோன்றுகிறது.

மிக முக்கியமாக, இது \emph{வெறும்} எண்ணத்தக்க தேர்வைத்
தாண்டி, \emph{முழுமையான} தேர்வை நியாயப்படுத்தாது. \emph{ஒவ்வொரு}
கணத்துக்கும் தேர்வுச் சார்பு வேண்டும் என்றால்,
``எந்த வரிசையெண் நீளத்துக்கும் நடைபெறும் அதிவேகப் பணி''
செய்ய முடியும் என்று கொள்ள வேண்டும். வெளிப்படையாகச்
சொன்னால், அந்தக் கருத்து நகைப்புக்குரியது.

\end{document}
